\documentclass{report}

\usepackage{hyperref}
\usepackage{appendix}

\title{System Requirements and Specifications \- \\ \large{MovieApp} }
\author{Alexander G. Ramirez}
\date{April 15th, 2026}

\begin{document}
\maketitle
\tableofcontents

\chapter{Abstract}


In order to show the process of documenting and transfering an API Specification
from \.NET MVC to another technology stack we will take a random API published on 
GitHub, document the functionality of the API, generate a suite of unit tests, 
and replicate the functionality in another technology.

For the API to be documented we will choose the 
\href{https://github.com/Iamkemical/MovieApp}{MoviaApp}\footnote{https://github.com/Iamkemical/MovieApp}
repository.  This is an old repository which contains many common \.NET MVC patterns.

For our specifications we will use the \href{https://www.openapis.org/}{OpenAPI} standard for docuemnting 
the API.  To document the high level (ie integrated) functionality we will use \href{https://cucumber.io/docs/}{Gherkin}.
Finally, we will also use UML 2.0 \href{https://en.wikipedia.org/wiki/Sequence_diagram}{Sequence Diagrams}.
This will serve to document the API interfaces (OpenAPI), the interaction and relationships between APIs 
(Sequence Diagrams), and the expected business or behaviour functionality (Gherkin) of the API.

Once the specifications and requirments are completed we will proceed to build implement the API in a new
programming language and framework.  In particular, we will develop the API using Rust.  Our unit tests 
will be developed in the Rust unit test framework.  For integration tests we will use the Gherkin features
and the Erlang Common Test framework.  Additionally, we will use the OpenAPI specs to automate test 
generation.

Once unit documentation is complete and unit tests are developed, we will proceed to create an Agile
backlog, prioritize functionality, and break development into user stories.

\chapter{Analysis}

\subsection{Overview}

Here we will analyze the general architecture of the API.

\chapter{Requirements}

\chapter{Tests}

\chapter{Implementation}

\end{document}